\begin{lemma}
Assume $k<g(k)$ and put $G(n)=g^{[n]}(k)$.  Every $l\in\mathbb N$ either
satisfies $l<G(0)$ or belongs to exactly one orbit block
$[G(n),G(n+1))$.
\end{lemma}
\begin{proof}
Fix $l\in\mathbb N$.  By \noderef{OrbitEmbedding}, choose an increasing
embedding $G$ such that

\[
G(j)=\mathsf{OrbitPoint}(g,k,j)=g^{[j]}(k)
\qquad\text{for every }j\in\mathbb N,
\]

where the second equality is the definition of \noderef{OrbitPoint}.
We prove the required disjunction by cases.

First suppose that $l<G(0)$.  This is the left side of the disjunction,
so the conclusion holds.

Now suppose that $G(0)\leq l$.  Applying \noderef{OrbitUnbounded} to
$g,k,hk$ and this $l$ gives an $r\in\mathbb N$ with

\[
l<\mathsf{OrbitPoint}(g,k,r)=G(r).
\]

The index $r$ cannot be $0$: if $r=0$, this inequality would say
$l<G(0)$, contradicting $G(0)\leq l$.  Hence $r=m+1$ for some
$m\in\mathbb N$, and therefore $l<G(m+1)$.  The set

\[
S=\{j\in\mathbb N\mid l<G(j+1)\}
\]

is consequently nonempty.  Let $n$ be its least element.  Then
$l<G(n+1)$.

We next prove $G(n)\leq l$.  If $n=0$, this is exactly the present
assumption $G(0)\leq l$.  If $n=t+1$, then $t<n$, so the minimality of
$n$ implies that $t\notin S$.  Thus $\neg(l<G(t+1))$, and linearity of
the order on $\mathbb N$ gives $G(t+1)\leq l$.  Since
$G(n)=G(t+1)$, we again obtain $G(n)\leq l$.

The two inequalities just proved are precisely the defining clauses of
the orbit block: by \noderef{OrbitBlock},

\[
\mathsf{OrbitBlock}(g,k,n,l)
\quad\Longleftrightarrow\quad
G(n)\leq l\ \text{ and }\ l<G(n+1).
\]

Thus at least one such index exists.  It remains to prove uniqueness.
Suppose that both $n$ and $m$ satisfy
$\mathsf{OrbitBlock}(g,k,\mathord\cdot,l)$.  Unpacking
\noderef{OrbitBlock} gives

\[
G(n)\leq l<G(n+1),
\qquad
G(m)\leq l<G(m+1).
\]

If $n<m$, then $n+1\leq m$.  Since $G$ is increasing by
\noderef{OrbitEmbedding}, this gives $G(n+1)\leq G(m)$.  Hence

\[
l<G(n+1)\leq G(m)\leq l,
\]

which is impossible.  If instead $m<n$, the same argument with the
indices interchanged gives

\[
l<G(m+1)\leq G(n)\leq l,
\]

again impossible.  Linearity therefore forces $n=m$, proving the
existence and uniqueness asserted in the right side of the disjunction.
\end{proof}
